\documentclass[11pt]{article}
\usepackage{amsmath,amssymb}
\title{The literal Wilson cubic vertex and its signed endpoint minor}\author{}\date{}
\begin{document}\maketitle
\paragraph{Original flat plaquette, no continuum replacement.}
Use a fixed quaternion product convention for
$SU(2)$. Let the four oriented fluctuation vectors in a plaquette
be $v_1,v_2,v_3,v_4$. Fixed links have zero fluctuation; inverse
occurrences include their original sign and insertion order.
The cubic coefficient of $1-\frac12\operatorname{ReTr}\prod_i
\exp(i t v_i\cdot\sigma)$, with the corresponding quaternion convention,
is
\[
 T_{3,p}=\sum_{i<j<k}(v_i\times v_j)\cdot v_k
       =F\cdot C,\quad
 F=\sum_i v_i,\quad C=\sum_{i<j}v_i\times v_j .
\]
A simultaneous reversal of the Lie-algebra convention reverses this
cubic sign but not any of the paired identities below.
Put
\[
 a=(v_1-v_3)/2,\quad b=(v_2-v_4)/2,\quad
 c=v_1+v_3,\quad d=v_2+v_4,\quad F=c+d .
\]
Direct expansion gives
\[
 C=2a\times b+(a+b)\times F+\tfrac12c\times d.
\]
Both last terms are orthogonal to $F$. Hence the ORIGINAL lattice
identity is
\[
 \boxed{T_{3,p}=2F\cdot(a\times b).}                              \tag{1}
\]
For a $\mu\nu$ plaquette this says exactly
\[
 F=\delta_\mu z_\nu-\delta_\nu z_\mu,\quad
 a=\operatorname{Av}_\nu z_\mu,\quad
 b=\operatorname{Av}_\mu z_\nu,
\]
with the two-point averages on the literal opposite links.
There is no higher-difference remainder. Zero extension of fixed
tangential links retains all boundary plaquettes. When a plaquette
has fewer than three free links the cubic polynomial vanishes by this
same identity; no boundary term is appended by hand.

\paragraph{All-cross Wick contraction before absolute values.}
At the flat covariant gauge with penalty one, the full relative vector
Green operator is diagonal in component and color. Consider the
$\delta_\mu z_\nu$ sector of (1) at plaquettes $x$ and $y$ of the same
orientation. Write $D=\delta_\mu$, $A=\operatorname{Av}_\mu$ on the
$\nu$ component grid. The unique $\mu$-component leg at each vertex
is averaged in direction $\nu$. Its kernel is denoted $G_\mu^{\rm av}$.

The two ways to pair the remaining $\nu$ legs have opposite signs from
the antisymmetric color tensor. Summing ALL cross-vertex pairings gives
\[
 \boxed{24G_\mu^{\rm av}(x,y)
 \left[(D_xD_yG_\nu)(A_xA_yG_\nu)
             -(D_xA_yG_\nu)(A_xD_yG_\nu)\right].}                \tag{2}
\]
The factor is $2^2$ from the two original vertices times six color
contractions. The two sectors of the curl have different component
multiplicities, so their cross pairing vanishes at this flat
component-diagonal reference. Different unordered plaquette
orientations likewise cannot match those component multiplicities.
This statement concerns the all-cross cubic graph; the completed
mean-current/tadpole, quartic, ghost and Haar terms remain separate
partners in the COMPLETE coefficient.

Let the $2$-by-$2$ matrix of endpoint Green values be
\[
 M=\begin{pmatrix}
 G_\nu(x,y)&G_\nu(x,y+\hat\mu)\\
 G_\nu(x+\hat\mu,y)&G_\nu(x+\hat\mu,y+\hat\mu)
 \end{pmatrix}.
\]
The two-point difference/average change has determinant minus one.
Therefore the bracket in (2) is EXACTLY
\[
 (D_xD_yG)(A_xA_yG)-(D_xA_yG)(A_xD_yG)=\det M.                    \tag{3}
\]
In particular the derivative-free and singly differentiated pieces
cancel algebraically before spatial bounds. This is an endpoint
minor, not automatically the fully contracted scalar $\Gamma$ of
the derivative network contact and boundary defect calculation. Treating it as that scalar energy without the remaining
index sums would be an unsupported step.

\paragraph{What changes at the actual constant-flux background.}
For one original background plaquette, move its background factors to
the right while retaining every prefix transport and the correct
before/after inverse insertion. The product becomes
\[
 \left[\prod_{i=1}^4\exp(i t\,\widetilde v_i\cdot\sigma)\right]
                    U_p^*,\qquad U_p^*=\exp(i\epsilon f_p\sigma_3).
\]
Each $\widetilde v_i$ is the ORIGINAL signed fluctuation parallel
transported along a prefix of at most four local links.
Let $V_3(\widetilde v)$ be the vector part of the coefficient of $t^3$
in the bracketed product. Direct scalar-vector multiplication gives
the exact local decomposition
\[
 T_{3,p,\epsilon}
 =2\cos(\epsilon f_p)\,
       \widetilde F\cdot(\widetilde a\times\widetilde b)
       +\sin(\epsilon f_p)\,[V_3(\widetilde v)]_3,                 \tag{4}
\]
with the consistent quaternion orientation.
Here $|V_3|\le C(\sum_i|\widetilde v_i|)^3$.
For the actual flux family, $f_p=R^{-2}$ on 01 plaquettes and zero
otherwise. Formula (4) retains a local explicit curvature vertex
alongside the transported curl vertex. It does not replace the
prefix transports by identity. Their background derivatives are
actual endpoint-current partners.

These prefix paths belong to the bounded local connector class of
the mixed grid joint current with connectors calculation. Thus the geometric scalar currents from this local transport
have a proved mixed-grid comparison. But the spin part of the
propagator, background derivatives of (4), changing ghost/projection
vertices, and the signed endpoint minor must still be combined.
The fact that the explicit sine amplitude has one factor $R^{-2}$
does not authorize dropping its mixed transport derivatives.

\paragraph{Verified cancellation and next quantitative join.}
The finite calculation
checks 200 exact original plaquette identities (1), the complete
color/component Wick factor (2), and 100 exact endpoint-minor identities.
These are algebraic regression checks, not a flux-response bound.
The unvaried graph now visibly contains two genuine derivative
differences. If the required current-weighted finite-boundary estimates
for their SIGNED combination were proved, this is the literal vertex
to which they would apply; derivatives may not simply be assigned to
the longest proper time.

The remaining main join is either a summation of (2)--(4) and the
ghost/spin contacts into the exact scalar energy cancellation of
the derivative network contact and boundary defect calculation plus controlled remainders, or a direct current-weighted
endpoint-minor estimate retaining those partners. The complete
boundary mass Ward and longitudinal cancellation remain in force.
No size-uniform bound on $\mathcal C_R''$, on the original remainder,
or on S94 is claimed by the vertex grouping alone.
\end{document}
